\newpage
\section{Otimização do Fokker 100}

O objetivo desta seção é minimizar o MTOW do Fokker 100 variando o
alongamento, $AR_w$, e a área de asa, $S_w$, sujeito a uma restrição de
envergadura, $b_w \le 30$~m, que é função das duas variáveis. O ponto de
partida é $AR_w = 7{,}5$ e $S_w = 90$~m$^2$. Formalmente, o problema de
otimização é definido como:

\begin{equation*}
\begin{aligned}
\underset{\vec{x}}{\text{minimizar}} \quad & \mathrm{MTOW}(\vec{x}),
    \qquad \vec{x} = \begin{bmatrix} AR_w & S_w \end{bmatrix}^{T} \\[1ex]
\text{sujeito a} \quad & g_1(\vec{x}) = \frac{b_w(\vec{x})}{30\,\mathrm{m}} - 1 \le 0 \\[1ex]
 & 7 \le AR_w \le 12 \\
 & 80\,\mathrm{m^2} \le S_w \le 120\,\mathrm{m^2}
\end{aligned}
\end{equation*}

Para melhorar o condicionamento numérico, as variáveis de projeto foram
normalizadas pelos valores do ponto de partida e a função objetivo pelo
MTOW inicial; a restrição $g_1$ já é adimensional por construção. A
otimização foi realizada com a função \texttt{minimize} da biblioteca
\textit{SciPy}, algoritmo SLSQP, com gradientes estimados por diferenças
finitas e parâmetros $\texttt{maxiter} = 200$ e $\texttt{ftol} = 10^{-8}$.
O histórico da otimização é apresentado na
Figura~\ref{fig:fokker100hist}.

\begin{figure}[H]
    \centering
    \includegraphics[width=0.85\textwidth]{images/fokker100_historico.png}
    \caption{Histórico da otimização ao longo das avaliações da função:
             variáveis de projeto normalizadas pelo valor inicial,
             objetivo e restrição de envergadura.}
    \label{fig:fokker100hist}
\end{figure}

Observa-se que o MTOW cai rapidamente nas primeiras iterações e estabiliza
em seguida, com as variáveis de projeto seguramente afastadas dos
batentes. A restrição $g_1$ aproxima-se de zero ao longo do processo, mas
recua levemente nas iterações finais, terminando \emph{inativa}:
no ótimo, $b_w = 29{,}89$~m e $g_1 = -3{,}8\times10^{-3}$. Como nenhum
batente está ativo e a norma do gradiente do objetivo no ponto final é
$\|\nabla f\| \approx 7\times10^{-6}$, conclui-se que \textbf{o ótimo não é
restringido}: trata-se de um mínimo interior, cuja condição de
otimalidade se reduz a $\nabla f = 0$, sem contribuição de
multiplicadores de Lagrange.

A Figura~\ref{fig:fokker100espaco} apresenta o espaço de projeto completo,
com o caminho percorrido pelo otimizador. O mínimo situa-se em uma bacia
rasa de MTOW, próxima --- porém aquém --- da fronteira $b_w = 30$~m, o que
confirma visualmente o caráter interior do ótimo e explica a convergência
em poucas iterações.

\begin{figure}[H]
    \centering
    \includegraphics[width=0.95\textwidth]{images/fokker100_espaco_projeto.png}
    \caption{Espaço de projeto $AR_w \times S_w$ colorido por MTOW, com a
             restrição de envergadura (vermelho) e a trajetória do
             otimizador do ponto de partida (verde) ao ótimo (azul).}
    \label{fig:fokker100espaco}
\end{figure}

Os resultados estão resumidos na Tabela~\ref{tab:otim_fokker100}. A maior
variação relativa ocorreu no alongamento ($+40\,\%$), enquanto a área de
asa foi levemente reduzida; o MTOW caiu $1\,187$~kgf.

\begin{table}[H]
\centering
\caption{Resultados da otimização do Fokker 100.}
\label{tab:otim_fokker100}
\renewcommand{\arraystretch}{1.25}
\begin{tabular}{l|cccc}
\hline
Ponto & $AR_w$ & $S_w$ [m$^2$] & MTOW [kgf] & $b_w$ [m] \\
\hline
Partida   &  7{,}50 & 90{,}00 & 43\,067 & 25{,}98 \\
Otimizado & 10{,}52 & 84{,}88 & 41\,880 & 29{,}89 \\
\hline
\end{tabular}
\end{table}

A melhoria relativa da função objetivo foi de

\begin{equation*}
    \frac{43\,067 - 41\,880}{43\,067} \times 100\,\% = 2{,}756\,\%.
\end{equation*}

O custo computacional da otimização foi de 33 avaliações da função
objetivo (\texttt{nfev}), correspondendo a 68 execuções da função
\texttt{analyze} (objetivo e restrição são avaliados em varreduras
separadas de diferenças finitas), em 11 iterações do SLSQP.
